\scanpage{233}
ce complexe est déduit d'une résolution flasque $\widetilde\Phi_*^{p*}$ de \deleted{\unsure{$\Phi$}} $k_*$ en \unsure{la|le} ``tronquant'' en \deleted{\unsure{dim|dimension}} degré $\leq p$, i.e. en prenant $\Phi_*^{p*}=\tau^{\leq p}\widetilde\Phi_*^{p*}$
\[
 \deleted{\unsure{\Phi^p}}\quad\Phi_*^{p\,i}=
 \begin{cases}\widetilde\Phi_*^{p,i}&\text{si }i<p\\
 Z^p(\widetilde\Phi_*^{p,*})&\text{si }i=p\\
 0&\deleted{\text{si}}\text{ si }i>p\end{cases}
\]

Par suite, pour tt $\unsure{\mathscr X|X|\mathcal X}$ semi-simplicial, on aura
\[
 \begin{aligned}
 \underset{\substack{\rotatebox{90}{$=$}\,\mathrm{défn}\\C^{p*}_{\mathrm{DRpd}}(\mathscr X,k)}}{\operatorname{Hom}(\mathscr X,\Phi_*^{p*})}
 &\simeq\deleted{\unsure{\mathcal C}}\operatorname{Hom}(X,\tau^{\leq\unsure{k|p}}\widetilde\Phi_*^{p*})
 \simeq\tau^{\leq p}\Gamma(X,\widetilde\Phi_*^{p*})\\
 &\simeq\tau^{\leq p}\mathbb R\unsure{\mathbb{\Gamma}|\Gamma}(\mathscr X,k)
 \end{aligned}
\]
\deleted{\unsure{que}} en particulier
\[
 H^{p*}_{\mathrm{DRpd}}(\mathscr X,k)\unsure{\text{---}|\longrightarrow}H^*(\mathscr X,k)
\]
\deleted{est un isom} tel que
\[
 H^{pi}_{\mathrm{DRpd}}(\mathscr X,k)\begin{cases}\xrightarrow{\sim} H^i(\mathscr X,k)&\text{si }i\leq p\\
 =0&\text{si }i>p.\end{cases}
\]

\begin{center}\rule{3cm}{0.4pt}\\[-0.7ex]\rule{1.6cm}{0.4pt}\end{center}

Soit $S$ anneau avec idéal $S^+$ à p.d, muni d'un \deleted{suite \unsure{décroissante|de générateurs}} $t\in S^+$. Quitte à remplacer $S$ par son séparé complété pour la top. engendrée par \unsure{les|ces} idéaux
\[
 S^{[i]}=\text{idéal engendré par les }t^{(j)}\quad(j\geq i)
\]
on va supposer $S$ séparé et complet pour la top définie par une suite décroissante d'idéaux $S^{[i]}$, chacun étant l'adhérence de l'idéal engendré par $t^{\unsure{(j)|[j]}}$ $(j\geq i)$.

% archive sheet number: 232

